\uuid{Feij}
\exo7id{839}
\titre{exo7 839}
\auteur{ridde}
\organisation{exo7}
\datecreate{1999-11-01}
\isIndication{false}
\isCorrection{false}
\chapitre{Calcul d'intégrales}
\sousChapitre{Primitives diverses}
\cpge{mpsi-mp2i.nombres-reels-suites.c, pcsi.nombres-reels-suites.b, ptsi.nombres-reels-suites.b, mpsi-mp2i.nombres-reels-suites.d, pcsi.nombres-reels-suites.c, ptsi.nombres-reels-suites.c, mpsi-mp2i.nombres-reels-suites.e, pcsi.nombres-reels-suites.d, ptsi.nombres-reels-suites.d, bcpst1.suites-reelles}
\module{Analyse}
\niveau{L1}
\difficulte{}

\contenu{
\texte{
Soit $a_{n} = \int_0^1 t^n e^t dt$.
}
\begin{enumerate}
    \item \question{Calculer $a_0, \ldots, a_4$.}
    \item \question{Etudier la suite $ (a_n)_{n \in \Nn}$.}
\end{enumerate}
}
